\section{Conclusion}
\label{sec:conclusion}

\note{This conclusion is provisional scaffold text. It should be rewritten after the empirical findings, contribution claims, and limitations are fully established.}

This paper is designed to study biometric voter registration in Brazil as both an administrative reform and a political event. The central premise is that changes to election infrastructure can affect not only the mechanics of registration, but also the social composition of the electorate and the way citizens think about institutions and democratic procedures. The project therefore links administrative data and survey evidence within a common framework organized around staggered municipal adoption.

The intended contribution is to keep together questions that are often studied separately. One strand of analysis focuses on electorate size and composition, while another focuses on trust, democratic support, and beliefs about electoral integrity. Bringing these strands into the same paper may clarify when administrative modernization appears primarily technical, when it appears politically consequential, and where the strongest inferential limits remain.

At this stage, the manuscript is a structured starting point rather than a finished argument. The next steps are to verify treatment timing, harmonize municipality identifiers, document data provenance, choose estimators transparently, and replace scaffold prose with evidence-backed claims. If the project succeeds, the final paper should offer a careful account of how biometric registration altered electoral administration in Brazil and what those changes may have meant for democratic confidence and skepticism.
